\rhead{RB-EL3: Swarm Intelligence \& Multi-Robot Systems}
\phantomsection 
\section*{Swarm Intelligence \& Multi-Robot Systems (RB-EL3)}
\label{elec:RB_3}

\noindent \small{\hyperref[main:scheme]{$\leftarrow$ Back to Scheme}} \vspace{0.5cm}

\noindent \textbf{Credits:} 4 (3-0-2)

\subsection*{Theory}
\noindent\textbf{Unit 1} \\
\textit{Swarm Intelligence Foundations:} Biological inspiration (ants, bees, birds, fish schooling), Boids model (separation, alignment, cohesion), Reynolds rules and flocking behavior, Particle Swarm Optimization (PSO: velocity update, inertia weight), Ant Colony Optimization (ACO: pheromone trails, evaporation), Artificial Bee Colony (ABC), Swarm stability analysis.

\vspace{0.3cm}

\noindent\textbf{Unit 2} \\
\textit{Multi-Robot Communication and Coordination:} Communication topologies (centralized, decentralized, hybrid), Consensus algorithms (average, leader election), Task allocation (CBBA, auction-based), Formation control (virtual structures, leader-follower), Potential field methods (artificial potentials, collision avoidance), Graph-theoretic approaches (Riemannian manifolds).

\vspace{0.3cm}

\noindent\textbf{Unit 3} \\
\textit{Swarm Robotics Algorithms:} Coverage control (Voronoi partitions, Lloyd algorithm), Foraging and patrolling strategies, Flocking with obstacles, Self-assembly and shape formation, Density control and deployment, Scalable coordination (local interactions, emergent behavior), Scalability analysis (communication overhead, convergence time).

\vspace{0.3cm}

\noindent\textbf{Unit 4} \\
\textit{Distributed Estimation and Mapping:} Multi-robot SLAM (centralized, decentralized, DDF-SLAM), Data association (JPDDA, Murty's algorithm), Cooperative localization (EKF/UKF fusion), Relative pose estimation (anchor nodes), Graph SLAM optimization (pose graph, factor graphs), Outlier rejection and robust estimation.

\vspace{0.3cm}

\noindent\textbf{Unit 5} \\
\textit{Advanced Topics and Applications:} Heterogeneity in swarms (specialized roles), Fault tolerance and resilience, Learning in multi-agent systems (MARL, QMIX), Swarm reinforcement learning, Real-world applications (search-rescue, agriculture, warehouse automation), Hardware platforms (Khepera, Kilobot, TurtleBot swarm), Standardization (ROS2 multi-robot, MQTT/IoT protocols).

\newpage

\subsection*{Practical}
\noindent\textbf{Lab 1: Boids Flocking Simulation} \\
Focuses on Unity/Gazebo implementation, parameter tuning.

\vspace{0.2cm}

\noindent\textbf{Lab 2: Particle Swarm Optimization} \\
Focuses on PyTorch/DEAP for robotics path planning.

\vspace{0.2cm}

\noindent\textbf{Lab 3: Ant Colony Path Finding} \\
Focuses on graph-based ACO for multi-robot navigation.

\vspace{0.2cm}

\noindent\textbf{Lab 4: Formation Control} \\
Focuses on leader-follower, virtual structure simulations.

\vspace{0.2cm}

\noindent\textbf{Lab 5: Task Allocation (CBBA)} \\
Focuses on auction algorithm implementation, conflict resolution.

\vspace{0.2cm}

\noindent\textbf{Lab 6: Multi-Robot Coverage} \\
Focuses on Voronoi-based area partitioning in Gazebo.

\vspace{0.2cm}

\noindent\textbf{Lab 7: Consensus and Synchronization} \\
Focuses on average consensus, phase synchronization.

\vspace{0.2cm}

\noindent\textbf{Lab 8: Cooperative SLAM} \\
Focuses on ROS2 multi-robot mapping, loop closure.

\vspace{0.2cm}

\noindent\textbf{Lab 9: Swarm Hardware Experiment} \\
Focuses on TurtleBot3 swarm, real-world flocking.

\vspace{0.2cm}

\noindent\textbf{Lab 10: Complete Swarm Application} \\
Focuses on search-and-rescue scenario with 10+ robots.

\vspace{0.2cm}

\newpage
\rhead{Curriculum Schema}
